% 皆さんは，スライディングブロックパズルを知っているだろうか．私は知っている．
% 有名なところで言えば「15パズル」だろう．
% \(4 \times 4\)のボード上に\(4 \times 4-1\)すなわち15枚の駒があり，1駒ぶんの空きを利用して駒をスライドし，駒を目的の配置にするものである．
% 少年時代，私は自分でそれを動かしながらもバラバラだった駒が整列していく様に感極まっていたことをよく覚えている．あれから10数年の時を経て，研究という大義の下に
% このパズルに帰着できたことは誠に光栄なことであり，今も興奮を抑えることができず手が震えているなんて恥ずかしくて言えない，言えないよそんなこと．

% 余談がすぎてしまった．本研究は，既に完成されている絵のスライディングブロックパズルを並べ替えることにより，別の新しい絵を描くものである．並び替える問題を割当問題に定式化し解を得て，そのパズルを解く様子をProcessingを用いて可視化した．
% さぁ今宵，貴方をアカデミックでファニーでインタレスティングでファンタジカルな私の研究へとお連れする．
皆さんは，スライディングブロックパズルを知っているだろうか．
有名なところで言えば「15パズル」だろう．
\(4 \times 4\)のボード上に\(4 \times 4-1\)すなわち15枚の駒があり，1駒ぶんの空きを利用して駒をスライドし，駒を目的の配置にするものである．
少年時代，私は自分でそれを動かしながらもバラバラだった駒が整列していく様に感極まっていたことをよく覚えている．

本研究は，既に完成されている絵のスライディングブロックパズルを並べ替えることにより，別の新しい絵を描くものである．並び替える問題は割当問題に定式化し解を得て，そのパズルを解く様子はProcessingを用いて可視化した．
学術的な要素だけでなく芸術面をも考慮された本研究を，是非楽しんでご覧になっていただきたい．